% =============================================================
% INTRODUCCIÓN — formato IEEE
% Para usar con \usepackage[backend=biber,style=ieee]{biblatex}
% y \addbibresource{referencias.bib}
% =============================================================

La evaluación constituye uno de los pilares fundamentales del proceso educativo, en tanto opera como mecanismo de regulación que articula la enseñanza con el aprendizaje y orienta la toma de decisiones pedagógicas. En las últimas tres décadas, el paradigma tradicional centrado en la medición de contenidos memorísticos ha cedido terreno frente a un enfoque competencial que prioriza la capacidad del estudiante para movilizar conocimientos, habilidades y actitudes en situaciones auténticas \cite{black1998inside,oecd2019compass}. Este desplazamiento, promovido por organismos internacionales como la OCDE, la UNESCO, el Banco Mundial y el Consejo de Europa, ha reconfigurado los marcos curriculares contemporáneos hacia una arquitectura basada en competencias que persigue el desarrollo de aprendices autónomos, críticos y preparados para los retos de un entorno dinámico \cite{unesco2022contract,voinea2025competency}.

Pese a la solidez teórica de este paradigma, su implementación efectiva en el aula sigue presentando desafíos significativos. Los resultados del \textit{Programme for International Student Assessment} (PISA) 2022 revelan un panorama preocupante: en los países de la OCDE el rendimiento medio en matemáticas cayó 15~puntos entre 2018 y 2022, la mayor caída registrada históricamente \cite{oecd2023pisa}. En América Latina y el Caribe la situación es aún más crítica, ya que tres de cada cuatro estudiantes de quince años no alcanzan el nivel mínimo de competencia matemática, frente al 31\,\% observado como promedio en la OCDE \cite{bid2023pisa,unesco2024lac}. Esta brecha evidencia la urgencia de repensar no solo qué se enseña, sino también cómo se evalúa el progreso de los estudiantes desde las primeras etapas formativas, dado que las habilidades matemáticas adquiridas en la educación básica predicen significativamente el desempeño académico futuro.

La investigación educativa contemporánea ha demostrado consistentemente que la evaluación formativa, entendida como un proceso cíclico de recopilación de evidencias y retroalimentación oportuna, produce ganancias sustanciales en el aprendizaje. Black y Wiliam \cite{black1998inside}, tras analizar más de 250 estudios empíricos, concluyeron que las prácticas de evaluación formativa de calidad pueden duplicar la velocidad del aprendizaje. Una revisión sistemática reciente publicada en \textit{ZDM~--~Mathematics Education} \cite{maskos2025fa}, que examinó 45 estudios sobre evaluación formativa en matemáticas durante el período 2015--2023, ratifica este efecto y destaca que las evaluaciones basadas en computador, en entornos de aprendizaje interactivos, presentan un potencial particularmente prometedor cuando se implementan de forma regular durante al menos tres meses con instrumentos alineados al contenido específico.

En este contexto, la incorporación de tecnologías digitales se presenta como una vía estratégica para superar las limitaciones operativas de la evaluación tradicional. Las plataformas de aprendizaje basadas en juego (\textit{digital game-based learning}) han mostrado mejoras significativas en la motivación, la participación cognitiva y el rendimiento matemático de estudiantes de educación primaria, en comparación con los métodos no digitales \cite{debrenti2024gbl,sayed2023adaptive}. De manera complementaria, la integración de inteligencia artificial generativa en los procesos de evaluación formativa abre nuevas posibilidades para entregar retroalimentación personalizada, adaptativa y oportuna, lo cual constituye uno de los factores más influyentes en el aprendizaje según el meta-análisis de Hattie y Timperley \cite{hattie2007feedback}. Estudios recientes confirman que los sistemas de IA aplicados a la evaluación pueden enriquecer el diseño de tareas, fomentar la metacognición del estudiante y fortalecer la literacidad evaluativa del docente \cite{maksimchuk2025genai,gonzalez2021ai}.

No obstante, una revisión cuidadosa del estado del arte revela que la mayoría de las herramientas tecnológicas existentes para evaluación por competencias~---~como SCALA, CAT, SECAT, MCAT, LACAMOLC y AEEA~---~han sido diseñadas predominantemente para la educación superior, dejando un vacío instrumental significativo en la educación básica elemental. Adicionalmente, en este nivel formativo concurren factores cognitivos y emocionales específicos~---~como la ansiedad matemática, la carga cognitiva elevada y las barreras de lectoescritura~---~que demandan soluciones especialmente diseñadas. La literatura sobre accesibilidad cognitiva y aprendizaje basado en imágenes ha demostrado que el uso de pictogramas y apoyos visuales reduce la carga cognitiva, facilita la asociación entre representaciones simbólicas y elementos del mundo real, y mejora significativamente la comprensión en niños, especialmente aquellos con necesidades educativas especiales \cite{mayer2014multimedia,guillomia2021cognitive}.

Frente a este escenario, el presente proyecto plantea el desarrollo de \textbf{Kiddy Quiz}, una aplicación web inclusiva orientada al control y evaluación de competencias curriculares en matemáticas para estudiantes de educación básica elemental. La propuesta se sustenta en cuatro pilares conceptuales: la estructuración de las evaluaciones por módulos y niveles de dificultad alineados con el currículo oficial, la integración de múltiples tipos de preguntas que favorezcan la evaluación auténtica, la incorporación de pictogramas de la base ARASAAC para asegurar la accesibilidad cognitiva infantil, y la utilización de la API generativa Gemini~2.5~Flash para proveer retroalimentación motivacional automatizada al término de cada evaluación. Adicionalmente, el sistema ofrece al docente analíticas detalladas, mapas de competencias y recomendaciones pedagógicas asistidas por inteligencia artificial, herramientas que automatizan tareas que manualmente resultan inviables para un profesor con numerosos estudiantes.

El desarrollo de Kiddy Quiz se gestionó bajo la metodología ágil \textit{Extreme Programming} (XP), priorizando ciclos cortos de retroalimentación con docentes expertos y la entrega iterativa de incrementos funcionales validados. La arquitectura técnica integra Angular para el \textit{frontend}, NestJS para la lógica del \textit{backend} y PostgreSQL como motor de base de datos relacional, complementada con servicios externos de Cloudinary para gestión multimedia y Google Gemini para procesamiento generativo. La protección de los datos se garantiza mediante autenticación con tokens JSON Web Token (JWT) y control de acceso basado en roles. La validación empírica del sistema se realizó a través de la Escala de Usabilidad del Sistema (SUS) y el Modelo de Aceptación Tecnológica (TAM) aplicados a docentes y estudiantes de educación básica elemental, instrumentos ampliamente reconocidos en la investigación en interacción humano-computador \cite{brooke1996sus,davis1989tam}.

El presente documento se estructura en cuatro capítulos. El Capítulo~I expone el marco contextual del proyecto, detallando el problema de investigación, los objetivos general y específicos, así como la justificación. El Capítulo~II desarrolla el marco teórico y referencial, abordando la evolución de las evaluaciones educativas, el enfoque por competencias en matemáticas, los tipos de preguntas y el uso de elementos multimedia, la tecnología educativa y la metodología XP. El Capítulo~III describe la metodología aplicada, el cronograma, el presupuesto, los instrumentos de seguimiento y los indicadores de evaluación. Finalmente, el Capítulo~IV presenta los resultados obtenidos, incluyendo las historias de usuario, los casos de uso extendidos, la arquitectura del sistema, el diseño del prototipo, la integración de Gemini~AI, los mecanismos de seguridad y los resultados de las pruebas de usabilidad y aceptación tecnológica con usuarios reales.